\documentclass{discussion}

\usepackage{amsmath, amsthm, amssymb}
% \usepackage{xcolor, hyperref, cleveref}
\usepackage{mathtools}
\usepackage{derivative}
\usepackage{tabularx}
\usepackage{graphicx}
\usepackage{tasks}
\usepackage{mdframed}
\usepackage{interval}
% \usepackage{tikz}
% \usetikzlibrary{patterns}
\usepackage{commands}

\settasks{
    % label-width = 13.9917pt,
    % item-indent = 17.64131pt,
    % after-item-skip = 80pt,
    % column-sep = 20pt,
    label = \boxed{\phantom{.}}
}

\discussion{10}
\author{Sections B04 and B06}
\date{Fall 2024}

\begin{document}

% \section{Lagrange Multipliers}

\begin{question}
    Consider \(f(x, y) = xy - 2x - y\) and \(g(x, y) = 2x + y\).
    \begin{subquestion}
        \item What is the maximum value of \(f(x, y)\) subject to the constraint \(g(x, y) = 20\)?
        \vspace{80pt}
        \item If \(f(x, y)\) is the profit as a function of minutes of advertising on TV (\(x\)) and podcast (\(y\)), and \(g(x, y)\) is the total cost of advertising (so \(xy\) is the revenue).
        What is the real-world interpretation of the Lagrange multiplier \(\lambda\)? (Select multiple and argue why they must be equal.)
        \begin{tasks}(2)
            \task The profit gained per minute of advertising on TV.
            \task The profit gained per minute of advertising on podcast.
            \task The profit gained per dollar spent on advertising on TV.
            \task The profit gained per dollar spent on advertising on podcast.
        \end{tasks}
    \end{subquestion}
\end{question}

\begin{question}
    Find the point on the graph of \(y = x^2 + 4x + 6\) that is the closest to the line \(y = 2x\).
\end{question}

\vspace{80pt}

\begin{question}
    Consider the tilted ellipse \(5x^2 - 6xy + 5y^2 = 4\).
    Find the points on the ellipse that are closest to and farthest from the origin.
\end{question}

\vspace{80pt}

% \section{Second derivative test}

\begin{question}[Second derivative test]
    Find the critical points of \(f(x, y) = x^2 + 3xy + y^2\) and determine whether they are local maxima, minima, or saddle points.
\end{question}

\clearpage

% \section{Review}

\begin{question}
    The following boxed arguments are incorrect or incomplete.
    Identify the error in each argument and provide a correct solution.

    \begin{mdframed}
        \textbf{Question:} Determine whether \(\sum_{n=2}^{\infty} \frac{1}{(\ln(n) + 2)^2}\) converges.

        \vspace{10pt}
        \noindent
        \textbf{Flawed Argument:} Yes, because as \(n\) approaches infinity, \(\ln(n) + 2\) approaches infinity, so \(\frac{1}{(\ln(n) + 2)^2}\) approach zero which is less than one.
    \end{mdframed}

    \vspace{10pt}

    \begin{mdframed}
        \textbf{Question:} Find the interval of convergence of the power series \(\sum_{n=0}^{\infty} \frac{(x-2)^n}{5n + 3}\).

        \vspace{10pt}
        \noindent
        \textbf{Flawed argument:} Apply ratio test, the series converges when
        \begin{equation*}
            \lim_{n\to \infty} \abs*{\frac{(x-2)^{n+1}}{5(n+1) + 3} \frac{5n + 3}{(x-2)^n}} = \lim_{n\to \infty} \abs*{\frac{5n + 3}{5n + 8} (x - 2)} = \abs{x-2} < 1.
        \end{equation*}
        Thus, \(-1 < x - 2 < 1\) which is equivalent to \(1 < x < 3\), so the interval of convergence is \(\ointerval{1}{3}\).
    \end{mdframed}

    \vspace{10pt}

    \begin{mdframed}
        \textbf{Question:} Calculate the limit \(\lim_{(x, y) \to (0,0)} \frac{xy}{x^2 + y^2}\) or show that it does not exist.

        \vspace{10pt}
        \noindent
        \textbf{Flawed argument:} Along the \(x\)-axis, the limit is \(\lim_{x \to 0} \frac{x \cdot 0}{x^2 + 0} = 0\) and along the \(y\)-axis, the limit is \(\lim_{y \to 0} \frac{0 \cdot y}{0 + y^2} = 0\), so the limit is \(0\).
    \end{mdframed}

    \vspace{10pt}

    \begin{mdframed}
        \textbf{Question:} Find the tangent line to the curve \(y = y^2 - x^3 + 1\) at the point \((1, 0)\).

        \vspace{10pt}
        \noindent
        \textbf{Flawed argument:} We compute the gradient of \(f(x, y) = y^2 - x^3 + 1\) which at the point \((1, 0)\) is
        \begin{equation*}
            \nabla f(x, y) |_{(1, 0)} = \veccomp{-3x^2, 2y}|_{(1,0)} = \veccomp{-3, 0}.
        \end{equation*}
        The tangent line is the line orthogonal to the gradient and passing through the point \((1, 0)\), so the equation of the line is \(\veccomp{x - 1, y - 0} \cdot \veccomp{-3, 0} = 0\) which is simply \(x = 1\).
    \end{mdframed}

    \vspace{10pt}

    \begin{mdframed}
        \textbf{Question:} Let \(f(x, y) = (x - 1)y\) and \(x(t) = \eu^{t}\) and \(y(t) = \cos(y)\).
        Find the first two nonzero terms of the Maclaurin series of \(f(x(t), y(t))\).

        \vspace{10pt}
        \noindent
        \textbf{Flaws argument:} We write \(x(t) = 1 + t\) and \(y(t) = 1 - \frac{t^2}{2}\), then we get
        \begin{equation*}
            f(x(t), y(t)) = (1 + t - 1)\left(1 - \frac{x^2}{2}\right) = t - \frac{t^3}{2}.
        \end{equation*}
    \end{mdframed}
\end{question}

% \begin{question}

% \end{question}

\end{document}